% Listings moved out of the main text by the fit pass (F3): rules of Sec. V-A.
\subsection{Channel Opening and Update}\label{lst:update}

These are the opening rule and the two rules of the update dance of
Sec.~\ref{sec:core-channels}.
\code{Lock_Funds_And_Open} creates the two mirrors \code{StateA} and \code{StateB} of state
\texttt{'0'}, the persistent facts that the other rules check, and the per-channel slot
\code{Mutex(ptr)} (Sec.~\ref{par:single-spend}).
\begin{spthy}
rule Lock_Funds_And_Open:
    [ ReadyToLock($A, $B, ptr, sA0, sB0, commitA0, commitB0),
      Wallet($A, ~fA), Wallet($B, ~fB) ]
    --[ ChannelOpen($A, $B, ptr), FundsLocked($A, $B, ptr),
        NewStateBuilt($A, $B, ptr, '0', ~fA, ~fB) ]->
    [ StateA($A, $B, ptr, '0', sA0, commitA0, ~fA, ~fB),
      StateB($A, $B, ptr, '0', sB0, commitB0, ~fA, ~fB),
      !ValidOnChain(ptr),
      !ChannelRecord(ptr, $A, $B),
      !ChannelConnect(ptr, $A, $B),
      !ChannelConnect(ptr, $B, $A),
      Mutex(ptr) ]
\end{spthy}
\code{Update_State} commits a new state: it samples the next state index and both fresh
revocation secrets, reads the two current mirrors and the new balances from the environment, and
produces two pending mirrors, each keeping the old state next to the new, signed one.
\begin{spthy}
rule Update_State:
    let newCommitA = <'commit', ptr, $A, $B, ~n_new, h(~sAnew)>
        newCommitB = <'commit', ptr, $B, $A, ~n_new, h(~sBnew)>
        sigAnew    = sign(newCommitA, ltkA)
        sigBnew    = sign(newCommitB, ltkB)
    in
    [ Fr(~n_new), Fr(~sAnew), Fr(~sBnew),
      !Ltk($A, ltkA), !Ltk($B, ltkB),
      StateA($A, $B, ptr, n_old, sAold, oldCommitA, fA, fB),
      StateB($A, $B, ptr, n_old, sBold, oldCommitB, fA, fB),
      In(<fA_new, fB_new>) ]
    --[ NewStateBuilt($A, $B, ptr, ~n_new, fA_new, fB_new) ]->
    [ PendingA($A, $B, ptr, n_old, sAold, oldCommitA,
               ~n_new, ~sAnew, newCommitA, sigAnew, fA_new, fB_new),
      PendingB($A, $B, ptr, n_old, sBold, oldCommitB,
               ~n_new, ~sBnew, newCommitB, sigBnew, fA_new, fB_new) ]
\end{spthy}
Only afterwards does \code{Revoke_Old_Secret} supersede the old state: it consumes the two
pending mirrors, installs them as the new state, and records, for both parties, the old
revocation secret as the persistent \code{!RevokedSecret} and the old commitment as the linear
token \code{RevokedCommitmentUnspent}. The signature checks of Sec.~\ref{sec:channel} become the
two \code{Eq} action facts: a \tamarin{} rule cannot fail, so restriction \code{Equality} turns
these action facts into hard prerequisites. The order (commit, then revoke) is what protects a
party that aborts mid-update: the old state is never revoked until a fully signed new one
exists.
\begin{spthy}
rule Revoke_Old_Secret:
    [ PendingA($A, $B, ptr, n_old, sAold, oldCommitA,
               n_new, sAnew, newCommitA, sigAnew, fA_new, fB_new),
      PendingB($A, $B, ptr, n_old, sBold, oldCommitB,
               n_new, sBnew, newCommitB, sigBnew, fA_new, fB_new),
      !Pk($A, pkA), !Pk($B, pkB) ]
    --[ StateUpdate($A, $B, ptr, n_new),
        Revoke($A, n_old, sAold),
        Revoke($B, n_old, sBold),
        Eq(verify(sigAnew, newCommitA, pkA), true),
        Eq(verify(sigBnew, newCommitB, pkB), true) ]->
    [ StateA($A, $B, ptr, n_new, sAnew, newCommitA, fA_new, fB_new),
      StateB($A, $B, ptr, n_new, sBnew, newCommitB, fA_new, fB_new),
      !RevokedSecret($A, n_old, sAold, oldCommitA),
      !RevokedSecret($B, n_old, sBold, oldCommitB),
      RevokedCommitmentUnspent(ptr, $A, $B, n_old, oldCommitA),
      RevokedCommitmentUnspent(ptr, $B, $A, n_old, oldCommitB) ]
\end{spthy}

\subsection{Channel Closing and Punishment}\label{lst:close}

The closing and punishment rules of Sec.~\ref{sec:core-channels}; the rules for party $B$
(\code{Publish_Current_B}, \code{Claim_Funds_B}, \code{Publish_Revoked_B}, \code{Punish_B}) are
symmetric. A unilateral close spends the current mirror: \code{Publish_Current_A} turns
\code{StateA} into an \code{OnChainDelayed} and an \code{OnChainInstant} output, and
\code{Claim_Funds_A} later releases the delayed output to its owner; its \code{CSVPassed} action
captures that time has passed.
\begin{spthy}
rule Publish_Current_A:
    [ StateA($A, $B, ptr, n, sA, commitA, fA, fB), !ValidOnChain(ptr) ]
    --[ Publish($A, n, commitA),
        PublishState($A, n),
        OnChainDelayedProduce($A, n),
        OnChainInstantProduce($B, n) ]->
    [ OnChainDelayed($A, n, commitA),
      OnChainInstant($B, n, commitA) ]
\end{spthy}
\begin{spthy}
rule Claim_Funds_A:
    [ OnChainDelayed($A, n, txContent) ]
    --[ ClaimFunds($A, n), CSVPassed($A, n), SettledTx($A, txContent) ]->
    [ FundsTo($A, n) ]
\end{spthy}
A cooperative close signs one shared settlement transaction instead.
\begin{spthy}
rule Honest_Close:
    let closetx = <'commit', ptr, $A, $B, n, fA, fB>
        signA   = sign(closetx, ltkA)
        signB   = sign(closetx, ltkB)
    in
    [ StateA($A, $B, ptr, n, sA, commitA, fA, fB),
      StateB($A, $B, ptr, n, sB, commitB, fA, fB),
      !ValidOnChain(ptr),
      !Ltk($A, ltkA), !Ltk($B, ltkB),
      !Pk($A, pkA),   !Pk($B, pkB) ]
    --[ Close($A, $B, ptr, n),
        OnChainInstantProduce($A, n),
        OnChainInstantProduce($B, n),
        Eq(verify(signA, closetx, pkA), true),
        Eq(verify(signB, closetx, pkB), true) ]->
    [ OnChainInstant($A, n, closetx),
      OnChainInstant($B, n, closetx) ]
\end{spthy}
Punishment publishes a revoked state, which consumes the linear token
\code{RevokedCommitmentUnspent} (the single-spend fix of Sec.~\ref{par:single-spend}: our first
version of \code{Publish_Revoked_A} had all of the premises below except this one), and then
transfers the whole channel to the honest party. Persistent knowledge of the revocation secret is
retained via \code{!RevokedSecret}; only spending the funds is linear.
\begin{spthy}
rule Publish_Revoked_A:
    [ !ChannelRecord(ptr, $A, $B),
      !ValidOnChain(ptr),
      !RevokedSecret($A, n_old, sAold, oldCommitA),
      RevokedCommitmentUnspent(ptr, $A, $B, n_old, oldCommitA) ]
    --[ PublishRevoked($A, n_old, oldCommitA),
        Cheat($A, n_old) ]->
    [ OnChainCheat($A, $B, n_old, oldCommitA),
      RevealedSecret($A, n_old, sAold, oldCommitA) ]
\end{spthy}
\begin{spthy}
rule Punish_A:
    [ OnChainCheat($A, $B, n_old, oldCommitA),
      RevealedSecret($A, n_old, sAold, oldCommitA) ]
    --[ Punish($B, $A, n_old) ]->
    [ AllFundsTo($B) ]
\end{spthy}

\subsection{Invoice, Settlement and Refund}\label{lst:settle}

The routing rules of Sec.~\ref{sec:core-routing} besides \code{Forward_HTLC} and
\code{AttackerSkip}, which are printed there. The invoice is tagged and names the receiver, the
sender and the payment hash; the sender checks the receiver's signature before it offers the
first HTLC.
\begin{spthy}
rule Receiver_Create_Invoice:
    let y = h(~x)
        invoice = <'invoice', $Receiver, $Sender, y>
        sigR    = sign(invoice, ltkR)
    in
    [ Fr(~x), !Ltk($Receiver, ltkR) ]
    --[ InvoiceCreated($Receiver, $Sender, y),
        SecretPreimage($Receiver, ~x) ]->
    [ ReceiverInv($Receiver, $Sender, ~x),
      Out(<'SecretRS', invoice, sigR>) ]

rule Sender_Offer_HTLC:
    let invoice = <'invoice', $Receiver, $Sender, y>
    in
    [ In(<'SecretRS', invoice, sigR>),
      !Pk($Receiver, pkR),
      !Ltk($Sender, ltkS),
      !ChannelConnect(ptrSF, $Sender, $Forward1),
      Mutex(ptrSF) ]
    --[ HTLCOffered($Sender, $Forward1, ptrSF, y),
        OfferFor($Sender, $Receiver, ptrSF, y),
        HTLCAdded(ptrSF, y),
        Neq($Sender, $Forward1),
        Neq($Sender, $Receiver),
        Neq($Forward1, $Receiver),
        Eq(verify(sigR, invoice, pkR), true) ]->
    [ SenderPending($Sender, $Forward1, $Receiver, ptrSF, y),
      Route($Sender, $Forward1, ptrSF, y),
      !RouteInfo($Sender, $Forward1, ptrSF, y) ]
\end{spthy}
The settle round runs backwards on \code{Fulfilled} facts. The receiver's and the intermediary's
settlement give the slot of the incoming channel back.
\begin{spthy}
rule Receiver_Fulfill_HTLC:
    [ Route($Prev, $Receiver, ptrIn, y),
      ReceiverInv($Receiver, $Sender, x) ]
    --[ HTLCFulfilled($Receiver, $Prev, ptrIn, y),
        PreimageReleased($Receiver, x),
        Redeemed(ptrIn, y),
        Old(ptrIn, y),
        Eq(h(x), y) ]->
    [ Fulfilled(ptrIn, x), Mutex(ptrIn) ]

rule Intermediary_Settle:
    [ IntermediaryPending($Me, $Prev, $Next, ptrIn, ptrOut, y, %vIn, %vOut, %fee),
      Fulfilled(ptrOut, x) ]
    --[ IntermediaryClaimed($Me, $Prev, ptrIn, x),
        UsedFulfilled(ptrOut, x),
        SettledHop(ptrIn, ptrOut, y),
        FeeEarned($Me, ptrIn, ptrOut, y, %fee),
        Redeemed(ptrIn, y),
        Old(ptrIn, y),
        FulfilledProduced(ptrIn, x),
        Eq(h(x), y) ]->
    [ Fulfilled(ptrIn, x), Mutex(ptrIn) ]

rule Sender_Settle:
    [ SenderPending($Sender, $Forward1, $Receiver, ptrSF, y),
      Fulfilled(ptrSF, x) ]
    --[ SenderSettled($Sender, $Forward1, $Receiver, ptrSF, x),
        UsedFulfilled(ptrSF, x),
        Eq(h(x), y) ]->
    [ ]
\end{spthy}
A timeout is an untimed event; after it, the pending HTLC on that channel is refunded to the
party that offered it, and the slot is given back.
\begin{spthy}
rule Open_Timeout:
    [ !ChannelConnect(ptr, A, B) ] --[ TimedOut(ptr) ]-> [ !Timeout(ptr) ]

rule Sender_Timeout_Refund:
    [ SenderPending($Sender, $Forward1, $Receiver, ptrSF, y), !Timeout(ptrSF) ]
    --[ Refunded($Sender, ptrSF, y),
        Old(ptrSF, y) ]->
    [ Mutex(ptrSF) ]

rule Intermediary_Timeout_Refund:
    [ IntermediaryPending($Me, $Prev, $Next, ptrIn, ptrOut, y, %vIn, %vOut, %fee),
      !Timeout(ptrOut) ]
    --[ Refunded($Me, ptrOut, y),
        RefundHop(ptrIn, ptrOut, y),
        Old(ptrOut, y) ]->
    [ Mutex(ptrOut) ]
\end{spthy}
